% \section{Conclusion}
% \label{sec:conclusion}

% We introduced \rcb{}, an interactive benchmark for evaluating whether coding agents can ask the right questions and build repositories from ambiguous product requirements. Unlike static from-scratch generation benchmarks, \rcb{} begins with fuzzy PRDs, withholds implementation details, and allows agents to clarify hidden constraints through a reproducible User Agent. Its 480 tasks span 12 programming languages, and its black-box test reconstruction enables language-agnostic functional evaluation.

% \rcb{} is designed to make ambiguity measurable. By combining public, native refined, and enhanced test suites with structural, semantic, design, and interaction-quality metrics, the benchmark exposes not only whether a repository works, but how the agent arrived there. This distinction is central to evaluating next-generation coding agents: robust software development requires agents to ask precise questions, remember clarified constraints, plan across long horizons, and translate incomplete intent into maintainable code. In this sense, interaction is not merely a user-interface convenience; it is a core software-engineering capability.

% Future work will expand the benchmark with richer human-in-the-loop interaction, more fine-grained ambiguity taxonomies, and longitudinal evaluations of agents that learn from failed trajectories.


\section{Conclusion and Future Work}
\label{sec:conclusion}

We introduced \rcb{}, an interactive benchmark evaluating how effectively coding agents build repositories from ambiguous product requirements. By using fuzzy PRDs and the reproducible User Agent protocol described in Section~\ref{sec:user-agent}, \rcb{} moves beyond static generation. Across 480 multi-language tasks, it measures not just whether the final code works, but whether agents can recover missing requirements and carry them through to a correct final artifact.

Future work will expand \rcb{} into end-to-end, multi-modal real-world scenarios. Specifically, by integrating with emerging frameworks like OpenClaw~\cite{openclaw2026}, we aim to introduce frontend development tasks and complex visual inputs, enabling richer human-in-the-loop evaluations for next-generation interactive agents.
